\chapter{Method}
\label{ch:method}

\section{Overview}
\label{sec:3.1}

\subsection{Methodological Objective}
\label{subsec:3.1.1}

The design has one goal: observe \ac{DNS} responses for a set of popular domains from geographically distributed probes, daily for 50 days. The credit budget and reproducibility constraints are both hard limits. Four empirical questions structure the work:

\begin{itemize}
  \item \textbf{Q1}: What proportion of Tranco top-ranked domains return geographically differentiated \ac{DNS} responses, and what infrastructure mechanisms explain the differentiation?
  \item \textbf{Q2}: What is the temporal stability of \ac{DNS} responses for popular domains at day, week, and month timescales?
  \item \textbf{Q3}: Do the geographic biases of the RIPE Atlas probe distribution significantly limit the ability to observe geographic \ac{DNS} variation?
  \item \textbf{Q4}: What is the measurable impact of resolver type --- local \ac{ISP} resolver versus public \ac{DNS} service --- on the \ac{DNS} responses observed?
\end{itemize}

\subsection{System Architecture}
\label{subsec:3.1.2}

The pipeline has five sequential stages, modelled loosely on OpenINTEL~\cite{vanRijswijk2016} but adapted to RIPE Atlas's credit-based model:

\begin{enumerate}
  \item \textbf{Domain corpus preparation}: Retrieval and filtering of the Tranco Top 250 domain list.
  \item \textbf{Measurement scheduling}: Configuration and deployment of RIPE Atlas \ac{DNS} measurement campaigns.
  \item \textbf{Data collection}: Automated daily retrieval of measurement results via the RIPE Atlas REST API.
  \item \textbf{Storage and preprocessing}: Parsing, validation, and archival in Parquet format.
  \item \textbf{Analysis}: Quantitative analysis addressing Q1--Q4.
\end{enumerate}

The RIPE Atlas credit budget caps total measurement volume. Adding domains means reducing probes or reducing frequency. Pick two of three: more domains, more probes, more measurements. The system is fully configurable via \texttt{--max-domains} and \texttt{--total-probes} parameters; the main phase operating point is 250 domains and 200 probes, consuming approximately 500,000 credits per day (250 domains $\times$ 200 probes $\times$ 10 credits per measurement).

\begin{table}[H]
\centering
\small
\begin{tabular}{c l l}
\toprule
\textbf{Step} & \textbf{Stage} & \textbf{Output} \\
\midrule
1 & Domain corpus preparation (\texttt{fetch\_tranco.py})         & \texttt{tranco\_corpus.csv} \\
2 & Measurement scheduling (\texttt{create\_ripe\_measurements.py}) & \texttt{measurements.json} \\
3 & Data collection (\texttt{fetch\_ripe\_atlas.py})               & \texttt{data/raw/msm\_*.json} \\
4 & Parsing \& storage (\texttt{parse\_dns\_results.py})           & \texttt{dns\_results*.parquet} \\
5 & Analysis (\texttt{analyse\_dns.py})                            & Reports \& figures \\
\bottomrule
\end{tabular}
\caption{Five-stage pipeline architecture. Steps 1--2 are run once at campaign initialisation; steps 3--4 run daily; step 5 runs weekly.}
\label{tab:pipeline_stages}
\end{table}

\subsection{Infrastructure}
\label{subsec:3.1.3}

The pipeline runs in a Docker container on a Raspberry Pi 5 at home --- a low-cost ARM64 board drawing under 10W continuously. No cloud VM: the container manages the full 50-day measurement campaign without manual intervention.

The image is built on \texttt{python:3.11-slim-bookworm}, compiled natively for ARM64. It carries only the six libraries the pipeline actually needs --- dnspython, pandas, pyarrow, scipy, matplotlib, and rclone --- and stays under 1~GB.

Results land in Docker named volumes, which survive image updates. rclone pushes a daily copy to Google Drive, so a Pi failure cannot lose more than 24 hours of data.

The pipeline is orchestrated by a cron daemon running inside the container, with the following schedule (all times UTC):

\begin{table}[H]
\centering
\small
\begin{tabular}{l l l}
\toprule
\textbf{Time (UTC)} & \textbf{Day} & \textbf{Action} \\
\midrule
10:00 & Daily & Fetch RIPE Atlas results + Parquet parsing \\
11:00 & Daily & Cloud sync (Parquet + raw JSON $\rightarrow$ Google Drive) \\
12:00 & Sunday & Q1--Q4 analysis + figure generation \\
13:00 & Sunday & Cloud sync (reports + figures $\rightarrow$ Google Drive) \\
04:00 & 1st of month & Log rotation \\
\bottomrule
\end{tabular}
\end{table}

The weekly Tranco list refresh is intentionally disabled during the active measurement period. Keeping the domain corpus fixed throughout the campaign is required for the validity of the Q2 temporal stability analysis: if domains entered or left the measurement set mid-campaign, observed response changes could not be distinguished from corpus changes. The corpus will be refreshed after the measurement period concludes.

The pipeline source code is available at \url{https://github.com/ogautier-unam/dns-measurements}. The five core scripts (\texttt{fetch\_tranco.py}, \texttt{create\_ripe\_measurements.py}, \texttt{fetch\_ripe\_atlas.py}, \texttt{parse\_dns\_results.py} and \texttt{analyse\_dns.py}) are introduced as they appear in the pipeline, alongside \texttt{pipeline.py}, which acts as the top-level orchestrator and is invoked by the cron scheduler.

\section{Domain Corpus Selection}
\label{sec:3.2}

\subsection{Tranco List Configuration}
\label{subsec:3.2.1}

The domain corpus is the Tranco Top 250 list, retrieved via the Tranco API at \url{https://tranco-list.eu}. The choice of Tranco over commercial ranking lists (Alexa --- discontinued May 2022 ---, Cisco Umbrella, Majestic) is motivated in Chapter 2 (Section~\ref{sec:2.6}) and rests on three properties established by Le Pochat~\cite{LePochat2019}: stability (0.6\% daily change versus approximately 50\% for post-2018 Alexa), manipulation resistance (quadrupled manipulation effort relative to single-source lists), and archivability (each generated list is assigned a unique identifier retrievable indefinitely via permalink).

\texttt{fetch\_tranco.py} downloads the most recent available list, with a seven-day fallback if that day's list is delayed. The corpus is then frozen for the entire campaign: mid-run domain turnover would conflate routing changes with corpus changes, making Q2 uninterpretable. The list identifier, generation date, and retrieval timestamp are written to \texttt{tranco\_meta.json}; any researcher can retrieve the exact same 250 domains from the Tranco permalink.

\subsection{Domain Corpus Size}
\label{subsec:3.2.2}

The empirical data collection was split into two sequential phases, tailored to the constraints of our RIPE Atlas credit allocations:

\textbf{Pilot phase (1--10 April 2026)}: To verify the end-to-end pipeline before committing significant resources, we restricted this initial run to the top 100 domains across 50 probes. Operating at this scale drew roughly 50,000 credits daily, allowing us to test probe selection, measurement initialisation, daily retrieval, parsing, and cloud synchronisation without exhausting our budget.

\textbf{Main phase (from 11 April 2026)}: With the pipeline validated and 30,000,000 RIPE Atlas credits secured, the corpus expanded to 250 domains across 200 probes --- 500,000 credits per day. 41 days at that rate costs 20,500,000 credits, leaving a 9,500,000-credit margin.

Why 250 domains? Calder~\cite{Calder2015} and Wang~\cite{Wang2018} show that \ac{CDN} routing diversity peaks at the top of the popularity ranking but measurable geographic variation extends into the broader popular tier. Top-250 captures the hyperscale operators --- Google, Cloudflare, Akamai, Meta --- while including second-tier platforms that geo-route with less pronounced differentiation. On the credit side: 250 $\times$ 200 probes $\times$ 10 credits = 500,000 per day; 41 days draws 20,500,000 of the 30,000,000 available, leaving a 9,500,000-credit reserve.

The pipeline exposes \texttt{--max-domains} and \texttt{--total-probes} parameters, so the corpus can be extended to a larger Tranco slice without any further modification.

\subsection{Domain Filtering and Validation}
\label{subsec:3.2.3}

Before the measurement campaign begins, each domain in the Tranco list is pre-validated by

\texttt{fetch\_tranco.py} in two steps:

\textbf{Technical validity}: A \ac{DNS} \ac{A} record query is issued against Google Public \ac{DNS} (8.8.8.8) and Cloudflare \ac{DNS} (1.1.1.1) as reference resolvers. Domains returning \texttt{NXDOMAIN} or \texttt{SERVFAIL} are excluded. Applied to the Tranco Top 10,000 list from which the campaign corpus is drawn, this step yielded 8,365 valid domains --- a filtering rate of approximately 16\%, higher than the 5--6\% reported by Le Pochat~\cite{LePochat2019}, which may reflect the broader scope of the list and the inclusion of parked or inactive domains in lower-ranked positions.

\ac{NS} records for each domain are resolved from the same two reference resolvers, mapping it to its authoritative server set. The result is saved as \texttt{tranco\_corpus.csv} with fields \texttt{rank}, \texttt{domain}, and \texttt{ns} (pipe-separated). Two of the 250 domains had authoritative \ac{NS} \ac{IP}s that failed to resolve at measurement time and were skipped automatically, leaving 248 active RIPE Atlas measurements.

\section{RIPE Atlas Measurement Configuration}
\label{sec:3.3}

\subsection{Probe Selection}
\label{subsec:3.3.1}

Probe selection is implemented in \texttt{create\_ripe\_measurements.py} and follows the filtering protocol recommended by Bajpai~\cite{Bajpai2017}. \textbf{System tag filtering.} All probes must satisfy the following RIPE Atlas system tags:

\begin{itemize}
  \item \texttt{system-ipv4-works}: \ac{IPv4} connectivity verified by built-in platform measurements.
  \item \texttt{system-resolves-a-correctly}: The probe's local resolver returns correct \ac{A} records, filtering out probes whose resolver intercepts or misroutes external queries~\cite{Boswell2024}.

\end{itemize}

Hardware probes must be version v3 or later. Under concurrent load, v1 and v2 hardware degrades by up to +7.7\,ms at the 95th percentile~\cite{Holterbach2015}; v3 probes show only +0.06\,ms. Software probes carry no such constraint.

\textbf{Geographic distribution.} To compensate for RIPE Atlas's documented geographic concentration (91\% of probes in European and North American networks~\cite{Bajpai2017}), probe allocation targets an inverse-weighted over-representation of under-represented regions. For the current campaign, distribution is delegated to RIPE Atlas using five built-in geographic zones:

\begin{table}[H]
\centering
\small
\begin{tabular}{l r}
\toprule
\textbf{RIPE Atlas zone} & \textbf{Allocation (200 probes)} \\
\midrule
West           & 60 (30\%) \\
North-Central  & 48 (24\%) \\
South-East     & 40 (20\%) \\
North-East     & 32 (16\%) \\
South-Central  & 20 (10\%) \\
\textbf{Total} & \textbf{200 (100\%)} \\
\bottomrule
\end{tabular}
\end{table}

\begin{figure}[H]
\centering
\includegraphics[width=0.85\linewidth]{figures/ripe_atlas_zones_map.png}
\caption{Geographic coverage of the five RIPE Atlas probe zones used for probe allocation. Source:~\cite{RIPEAtlasDocs2026}.}
\label{fig:ripe_atlas_zones}
\end{figure}

The resulting continent-level distribution is reported in Table~\ref{tab:4.2}. For future campaigns, the pipeline implements a continent-level stratification (Step 3) designed to achieve the following targets:

\begin{table}[H]
\centering
\small
\begin{tabular}{l l l l l}
\toprule
\textbf{Region} & \textbf{Atlas distribution} & \textbf{Pilot (50 probes)} & \textbf{Main (200 probes)} & \textbf{Target \%} \\
\midrule
Europe        & \textasciitilde{}40\% & 15  & 60  & 30\% \\
North America & \textasciitilde{}28\% & 12  & 50  & 25\% \\
Asia-Pacific  & \textasciitilde{}16\% & 10  & 40  & 20\% \\
South America & \textasciitilde{}8\%  &  5  & 20  & 10\% \\
Africa        & \textasciitilde{}5\%  &  5  & 20  & 10\% \\
Oceania       & \textasciitilde{}3\%  &  3  & 10  &  5\% \\
\textbf{Total} & \textbf{100\%} & \textbf{50} & \textbf{200} & \textbf{100\%} \\
\bottomrule
\end{tabular}
\end{table}

These proportions preserve the relative weighting from prior work~\cite{Bajpai2017} while scaling to the target probe count. The probe count is configurable via \texttt{--total-probes}.

\textbf{Country and \ac{AS} diversity.} Within each continental stratum, a three-round greedy algorithm maximises country count first, then \ac{AS} diversity:

\begin{itemize}
  \item \textbf{Round 1 --- one per country}: One probe per country, skipping micro-states (Luxembourg, Monaco, Gibraltar, etc.) whose \ac{DNS} infrastructure routes through the same \ac{CDN} \ac{PoP}s as their larger neighbours.
  \item \textbf{Round 2 --- micro-states}: One probe per previously skipped country, if the continental quota remains unfilled.
  \item \textbf{Round 3 --- fill with \ac{AS} diversity}: Additional probes per country, always preferring a new \ac{AS}, until the quota is met.
\end{itemize}

Within each round, the probe selected for a given country is the one with the highest hardware version that introduces a new \ac{AS} not yet represented in the selected set. The selected probe set is saved as \texttt{selected\_probes\_<date>.json} for full reproducibility. This algorithm is fully implemented in the pipeline and will be activated for future campaigns, replacing the zone-based assignment used for the current campaign.

\subsection{Primary Campaign: Direct Authoritative DNS Queries}
\label{subsec:3.3.2}

For Q1--Q3, queries go directly to the authoritative name server. The full trade-off is in Section~\ref{sec:2.5}. Briefly: local resolvers introduce two problems. First, centralisation~\cite{Xu2023}: 90\% of forwarding resolvers pass through 5\% of recursive resolvers, so two probes on opposite continents using local resolvers may reach the same upstream and get identical answers regardless of location. Second, cache state: differently-aged caches produce different answers not because of geographic routing, but because of timing artefacts. Direct authoritative queries skip both.

The default resolver approach is used for Q4 (Section~\ref{subsec:3.3.3}), where the research question is precisely about the impact of resolver choice on end-user experience.

For each domain in the corpus, the primary authoritative name server identified during pre-resolution is used as the measurement target. The RIPE Atlas \ac{DNS} measurement is configured as follows:


\begin{lstlisting}[style=json]
{
  "type": "dns",
  "af": 4,
  "target": "<authoritative_ns_ip>",
  "query_type": "A",
  "query_class": "IN",
  "query_argument": "<domain>",
  "use_probe_resolver": false,
  "set_rd_bit": false,
  "set_nsid_bit": true,
  "set_do_bit": false,
  "set_cd_bit": false,
  "protocol": "UDP",
  "udp_payload_size": 1024,
  "include_abuf": true,
  "retry": 2,
  "timeout": 5000,
  "is_oneoff": false,
  "interval": 86400,
  "description": "DNS geo-diversity - auth - <domain> - <YYYY-MM-DD>"
}

\end{lstlisting}

Table~\ref{tab:msm_params} describes each parameter and its justification.

Full API documentation is available in~\cite{RIPEAtlasDocs2026}.

\begin{table}[H]
\centering
\small
\begin{tabularx}{\linewidth}{l l X}
\toprule
\textbf{Parameter} & \textbf{Value} & \textbf{Justification} \\
\midrule
\texttt{type} & \texttt{"dns"} & RIPE Atlas measurement type for \ac{DNS} queries. \\
\texttt{af} & \texttt{4} & \ac{IPv4} address family; probes resolve the target using \ac{IPv4}. \\
\texttt{target} & \textit{ns\_ip} & \ac{IP} address of the authoritative name server identified during pre-resolution; queries bypass recursive resolvers entirely. \\
\texttt{query\_type} & \texttt{"A"} & Requests \ac{IPv4} address records, which are the primary indicator of \ac{CDN} geographic routing. \\
\texttt{query\_class} & \texttt{"IN"} & Internet class; the only class used in practice for public \ac{DNS}. \\
\texttt{query\_argument} & \textit{domain} & The domain name being queried. \\
\texttt{use\_probe\_resolver} & \texttt{false} & Disables routing through the probe's local resolver; queries are sent directly to the authoritative server. \\
\texttt{set\_rd\_bit} & \texttt{false} & \ac{RD} bit unset; the authoritative server is queried directly and must not recurse. \\
\texttt{set\_nsid\_bit} & \texttt{true} & Requests the \ac{NSID} option~\cite{RFC5001}; the authoritative server's response identifies the physical anycast instance that answered, enabling anycast routing analysis~\cite{Bortzmeyer2013,Finnegan2018}. \\
\texttt{set\_do\_bit} & \texttt{false} & \ac{DNSSEC} OK bit unset; \ac{DNSSEC} validation is not required for \ac{IP} address collection and would increase response size without benefit. \\
\texttt{set\_cd\_bit} & \texttt{false} & Checking Disabled bit unset; \ac{DNSSEC} checking is not relevant for direct authoritative queries. \\
\texttt{protocol} & \texttt{"UDP"} & Standard \ac{DNS} transport; UDP is used by default and is sufficient for responses of this size. \\
\texttt{udp\_payload\_size} & \texttt{1024} & EDNS0 buffer size in bytes; set large enough to accommodate the \ac{NSID} option in the response without truncation. \\
\texttt{include\_abuf} & \texttt{true} & Includes the raw \ac{DNS} response packet (base64-encoded), enabling full parsing of all \ac{DNS} sections and OPT record \ac{NSID} decoding~\cite{RFC5001}. \\
\texttt{retry} & \texttt{2} & Number of retransmission attempts on timeout before the result is recorded as no-response. \\
\texttt{timeout} & \texttt{5000} & Per-attempt timeout in milliseconds (5 seconds); consistent with RIPE Atlas defaults for \ac{DNS} measurements. \\
\texttt{is\_oneoff} & \texttt{false} & Configures the measurement as periodic rather than a one-shot query. \\
\texttt{interval} & \texttt{86400} & Repetition interval in seconds (24 hours); provides the daily temporal resolution needed to detect \ac{DNS} response changes within the credit budget. \\
\texttt{description} & \textit{see value} & Human-readable label in the RIPE Atlas dashboard, including the domain name and creation date for traceability. \\
\bottomrule
\end{tabularx}
\caption{RIPE Atlas DNS measurement parameters for the primary Q1--Q3 campaign.}
\label{tab:msm_params}
\end{table}

\subsection{Supplementary Campaign for Q4: Resolver Comparison}
\label{subsec:3.3.3}

To address Q4, a supplementary one-off measurement campaign is conducted on the full 250-domain corpus, identical to the primary campaign. Using the same domain set ensures full comparability between the authoritative \ac{DNS} responses (Q1--Q3) and the recursive resolver responses (Q4). This campaign provides four simultaneous resolver-type observations from the same 200 probes:

\begin{table}[H]
\centering
\small
\begin{tabular}{l l l l}
\toprule
\textbf{Campaign} & \textbf{Target} & \textbf{\texttt{use\_probe\_resolver}} & \textbf{Resolver} \\
\midrule
\texttt{isp\_local} & --- & \texttt{true} & Probe's local \ac{ISP} resolver \\
\texttt{google} & 8.8.8.8 & \texttt{false} & Google Public \ac{DNS} \\
\texttt{cloudflare} & 1.1.1.1 & \texttt{false} & Cloudflare \ac{DNS} \\
\texttt{quad9} & 9.9.9.9 & \texttt{false} & Quad9 \\
\bottomrule
\end{tabular}
\end{table}

The Q4 campaign uses \texttt{set\_rd\_bit: true} and a larger \ac{UDP} payload size (4096 bytes) appropriate for recursive resolver queries. Measurements are configured as one-off (\texttt{is\_oneoff: true}) to avoid continuous credit consumption.

The four configurations give a direct answer to Q4. Does resolver choice change which \ac{CDN}-assigned \ac{IP} a probe sees? Does the local \ac{ISP} resolver stay closer to the authoritative answer than Google or Cloudflare?

\subsection{Credit Budget}
\label{subsec:3.3.4}

The RIPE Atlas credit consumption is estimated as follows (10 credits per probe per measurement instance):

\begin{table}[H]
\centering
\small
\begin{tabular}{l l l l l}
\toprule
\textbf{Phase} & \textbf{Domains} & \textbf{Probes} & \textbf{Duration} & \textbf{Credits} \\
\midrule
Pilot --- primary (direct auth)       & 100 &  50 & 9 days  & $\approx$450,000 \\
Main --- primary (direct auth)        & 250 & 200 & 41 days  & $\approx$20,500,000 \\
Main --- Q4 resolver comparison       & 250 & 200 & once     & $\approx$2,000,000 \\
\textbf{Total campaign}             &     &     &          & $\approx$\textbf{22,950,000} \\
\bottomrule
\end{tabular}
\end{table}

The total estimated consumption of approximately 22,950,000 credits remains within the initial balance of 30,000,000 credits, leaving a margin of approximately 7,050,000 credits. The credit budget is monitored by the log output of \texttt{create\_ripe\_measurements.py}, which prints the estimated daily and total consumption before any measurements are created.

\section{Data Collection and Processing}
\label{sec:3.4}

\subsection{Daily Collection Pipeline}
\label{subsec:3.4.1}

Measurement results are collected from the RIPE Atlas REST API by \texttt{fetch\_ripe\_atlas.py} on a daily schedule at 10:00 UTC, retrieving results for the previous calendar day. The script reads the measurement plan from \texttt{measurements.json} (produced by \texttt{create\_ripe\_measurements.py}) and iterates over all active measurement identifiers.

The RIPE Atlas API paginates results; the script follows the \texttt{next} link chain until the day's window is exhausted. An \ac{HTTP} 429 pauses execution for the \texttt{Retry-After} duration. Any other failure gets five exponential-backoff retries; if all fail, the measurement ID goes into \texttt{fetch\_errors\_<date>.json}. Files already on disk are not re-fetched unless \texttt{--force} is set.

The raw results are stored in JSON files, which are kept in the data/raw/ folder and named \texttt{data/raw/m\-sm\_<id>\_<date>.json}.

\subsection{Parsing and Field Extraction}
\label{subsec:3.4.2}

Each raw RIPE Atlas \ac{DNS} result is parsed by \texttt{parse\_dns\_results.py}, which extracts the following fields into a structured Parquet record:

\begin{table}[H]
\centering
\small
\begin{tabular}{l l l}
\toprule
\textbf{Field} & \textbf{Type} & \textbf{Description} \\
\midrule
\texttt{msm\_id} & int64 & RIPE Atlas measurement ID \\
\texttt{prb\_id} & int32 & RIPE Atlas probe ID \\
\texttt{timestamp} & int64 & Actual measurement time (Unix epoch) \\
\texttt{date} & string & ISO date (YYYY-MM-DD) \\
\texttt{country\_code} & string & ISO 3166-1 alpha-2 country code \\
\texttt{asn\_v4} & int32 & Probe \ac{AS} Number \\
\texttt{continent} & string & Continent code (EU/NA/SA/AF/AP/OC) \\
\texttt{query\_domain} & string & Queried domain name \\
\texttt{rcode} & int16 & \ac{DNS} response code (0 = NOERROR) \\
\texttt{answer\_count} & int16 & Number of \ac{A}/\ac{AAAA} records returned \\
\texttt{answer\_ips} & string & Pipe-separated list of returned \ac{IP} addresses \\
\texttt{ttl\_min} & int32 & Minimum \ac{TTL} across answer records \\
\texttt{rt\_ms} & float32 & Response time in milliseconds \\
\texttt{nsid\_str} & string & \ac{NSID} value (ASCII, if returned) \\
\texttt{nsid\_hex} & string & \ac{NSID} value (hex encoding) \\
\texttt{abuf\_size} & int32 & \ac{DNS} response packet size in bytes \\
\texttt{resolver\_ip} & string & Target resolver \ac{IP} (Q4 campaign only) \\
\texttt{use\_probe\_resolver} & bool & True for \ac{ISP} local resolver campaign \\
\bottomrule
\end{tabular}
\end{table}

The \ac{NSID} value is extracted by decoding the raw ABUF binary field: the script walks the \ac{DNS} packet structure to locate the OPT record (type 41) in the Additional section, then parses EDNS0 options to extract option code 3 (\ac{NSID}, ~\cite{RFC5001}).

\subsection{Data Validation}
\label{subsec:3.4.3}

The following validation filters are applied during parsing:

\textbf{RCODE filtering}: Only results with \texttt{rcode = 0} (NOERROR) and \texttt{answer\_count > 0} are used for Q1 and Q2 content analysis. Results with \texttt{rcode != 0} are retained in the Parquet dataset with their specific error codes for reliability and censorship analysis.

\textbf{Deduplication}: Each daily append deduplicates on \texttt{(msm\_id, prb\_id, timestamp)}, where \texttt{timestamp} is the probe-side measurement clock, not the fetch clock. Re-fetching yesterday's results produces identical tuples; the deduplication step drops them without complaint.

\subsection{Storage Architecture}
\label{subsec:3.4.4}

Results are stored in a single-tier Parquet architecture optimised for the analytical query patterns of the analysis phase:

\textbf{Cumulative Parquet file} (\texttt{data/processed/dns\_results.parquet}): All parsed results are appended daily to a single Parquet file using pyarrow with a fixed schema. This file grows at approximately 40--60 MB/month for the primary campaign (250 domains $\times$ 200 probes $\times$ 30 days, after columnar compression with dictionary encoding for repeated string values such as \texttt{country\_code}, \texttt{continent}, and \texttt{answer\_ips}).

Daily per-date files (\texttt{dns\_results\_<date>.parquet}) sit alongside the cumulative file. A cumulative file corruption --- a write failure, a Docker volume issue --- is recoverable by replaying the daily appends in order.

Raw JSON files stay on the Pi for seven days. That window is long enough to re-parse yesterday's results if a schema fix is needed. After cloud sync confirms the copy on Google Drive, they are removed.

rclone pushes Parquet files, daily archives, and reports to Google Drive at 11:00 UTC. The credentials live in a named Docker volume and survive container image updates.

\textbf{Storage estimate} for a 50-day campaign (250 domains $\times$ 200 probes):

\begin{table}[H]
\centering
\small
\begin{tabular}{l l l}
\toprule
\textbf{Tier} & \textbf{Volume} & \textbf{Location} \\
\midrule
Raw JSON (7-day rolling) & \textasciitilde{}70 MB & Pi local \\
Cumulative Parquet & \textasciitilde{}80--100 MB & Pi + Google Drive \\
Daily Parquet archives & \textasciitilde{}80--100 MB & Pi + Google Drive \\
Reports and figures & \textasciitilde{}50 MB & Pi + Google Drive \\
\bottomrule
\end{tabular}
\end{table}

\bigskip
\section{Analysis Methods}
\label{sec:3.5}

\texttt{analyse\_dns.py} is step 5 in Table~\ref{tab:pipeline_stages}. It reads \texttt{dns\_results.parquet} and writes one CSV report and one or more PNG figures per question.

\subsection{Statistical Testing Approach}
\label{subsec:3.5.0}

Comparisons that pit two groups against each other --- probes from different continents, or different resolver types --- use the Mann-Whitney U test. It asks whether one group tends to produce larger values than the other, with no assumption about distributional shape.

Two reasons ruled out the t-test. The primary response variable, unique \ac{IP} count per domain, is right-skewed: most domains return one or two addresses, a handful of \ac{CDN} operators push the tail far right. And the continental groups differ substantially in size. Mann-Whitney makes no normality assumption and handles unequal groups.

The null hypothesis is identical response distributions in both groups. Significance threshold: $p < 0.05$.

\subsection{Q1 --- Geographic Diversity of DNS Responses}
\label{subsec:3.5.1}

Q1 asks whether the returned \ac{IP} address set shifts with probe geography. The analysis runs at two granularities drawn from the probe set in Section~\ref{subsec:3.3.1}.

At continent level, probes are grouped by region for each (domain, day); a domain is continent-diverse on that day if at least two continental groups return non-overlapping \ac{IP} sets. At country level, the same check runs for each of the 173 countries. Country-level analysis catches routing splits invisible at continent resolution --- a probe in France and a probe in Germany can return different \ac{IP}s while both label as ``Europe''.

Both granularities are then aggregated across time. For each domain, the fraction of measurement days showing geographic differentiation (\texttt{pct\_geo\_diverse}) measures how consistent the routing split is over the 50-day window. A single day of differentiation qualifies a domain as geographically diverse; the time series shows whether it is isolated or persistent.

For domains with geographic \ac{IP} variation, \ac{NSID} strings in the OPT record distinguish the mechanism. Varying \ac{NSID} per region means different anycast \ac{PoP}s were reached; uniform \ac{NSID} with different \ac{A} records points to a \ac{CDN} geo-database lookup from the same server.

Reported figures: proportion of domains showing continent-level and country-level diversity; proportion with intra-continental-only differentiation; distribution of distinct \ac{IP} set counts at both granularities; the top-20 most geographically diverse domains; and the \ac{NSID} rate split between diverse and uniform domains.

\subsection{Q2 --- Temporal Stability}
\label{subsec:3.5.2}

Q2 tracks change over time. For each (domain, probe) pair, consecutive daily observations are compared; a change event is recorded whenever the observed \texttt{answer\_ips} set differs from the previous measurement. The fraction of such pairs that differ is the per-probe daily change rate.

Three temporal windows are reported separately. Consecutive-day pairs measure short-cycle rotation. Seven-day gaps test whether whatever changed on day~1 is still changed a week later. Thirty-day gaps are the monthly stability estimate; at 50 days total, only 15 such pairs are available, but that is enough to distinguish stationary churn from a convergent process.

Change rates are then summarised at the domain level across the corpus. The distribution matters as much as the central estimate: a bimodal split between stable and volatile domains tells a different story than a uniform spread at the same mean.

Reported figures: the per-domain daily, weekly, and monthly change-rate distributions; proportion of fully stable domains (exactly 0\% across the window); corpus median; and the size of the zero-percent bin.

\subsection{Q3 --- Assessment of RIPE Atlas Geographic Bias}
\label{subsec:3.5.3}

Q3 asks whether the EU+NA concentration of RIPE Atlas probes distorts what the measurement sees. Two Mann-Whitney tests answer it at the same two granularities as Q1.

The first test compares per-domain unique \ac{IP} counts between EU+NA probe groups and all other regions combined (Asia-Pacific, Africa, South America, Oceania). The achieved deployment distribution is first reported by continent --- how far it deviates from the zone-based allocation target --- alongside per-continent coverage rates and mean unique-\ac{IP} counts. A significant result means the two halves of the probe set observe genuinely different \ac{DNS} responses; a non-significant result means probe imbalance does not bias the diversity estimate.

The second test uses a natural split in the dataset. The three-round selection algorithm (Section~\ref{subsec:3.3.1}) places one probe per country in Round~1 and adds extras in Round~3. The data therefore contains 34 single-probe countries and 139 countries with at least two probes. Round-3 extras tend to land in already-well-connected nations. The Mann-Whitney test on these two groups asks directly: does a country-first Round-1 probe add real geographic signal, or just a flag on the map?

Reported figures: actual probe distribution by continent and country; count of countries with at least one probe; single-probe versus multi-probe country counts per continent; mean unique \ac{IP}s per domain per continent; and both test statistics with p-values.

\subsection{Q4 --- Impact of Resolver Choice}
\label{subsec:3.5.4}

Q4 measures how much switching resolver type changes the answer a probe receives. The one-off campaign (Section~\ref{subsec:3.3.3}) runs the same 200 probes against each of the 250 domains four times --- once per resolver type (\texttt{isp\_local}, \texttt{google}, \texttt{cloudflare}, \texttt{quad9}).

Each resolver type is first profiled independently: error rate (non-zero RCODE fraction), median \ac{RTT}, and mean answer count. These baselines establish whether any resolver is structurally noisier or slower before the cross-resolver comparison begins.

The core comparison is pairwise. For each (domain, probe) pair the local \ac{ISP} \texttt{answer\_ips} is compared against each public-resolver response; the divergence rate is the fraction of pairs where the two sets differ. A geographic breakdown by continent answers whether the ISP-versus-public gap is uniform across regions or concentrated in specific areas, as would be expected if public resolvers steer probes to US-region \ac{CDN} nodes regardless of probe location.

A Mann-Whitney U test on \ac{RTT} distributions rounds out the comparison. Reported figures: error rate and median \ac{RTT} per resolver; pairwise \ac{IP} divergence rates; and the continental disaggregation of those rates.

\bigskip
\section{Ethical Considerations and Reproducibility}
\label{sec:3.6}

\subsection{Ethical Design}
\label{subsec:3.6.1}

The measurement design follows the ethical framework of Kisteleki~\cite{Kisteleki2016} and the principles of the Menlo Report (2012)~\cite{Bailey2012}:

All query targets are publicly accessible Tranco Top-250 domains. Direct authoritative queries for \ac{A} records are indistinguishable from normal client \ac{DNS} traffic; probe hosts are not running queries on behalf of the researcher that they would not otherwise run.

The campaign generates approximately 50,000 queries per day across the full corpus. OpenINTEL's production campaign runs several hundred million queries per day~\cite{vanRijswijk2016}; the impact on authoritative infrastructure is negligible.

Every measurement carries a project description in the RIPE Atlas dashboard. Authoritative server operators can identify the source from that tag. The full measurement plan is in \texttt{measurements.json}.

\subsection{Scientific Reproducibility}
\label{subsec:3.6.2}

The Tranco list identifier in \texttt{tranco\_meta.json} pins the domain corpus; any researcher can pull the exact same 250 domains from the Tranco permalink years later.

All RIPE Atlas measurement IDs are in \texttt{measurements.json}. Any researcher with a RIPE Atlas account can retrieve the original raw results using those IDs directly. The probe identifiers, country, \ac{AS} number, and hardware version for every selected probe are logged to \texttt{selected\_probes\_<date>.json}.

\texttt{requirements.txt} pins Python 3.11, dnspython, pandas, pyarrow, scipy, and matplotlib to specific versions --- no floating dependencies. The \texttt{Dockerfile} builds for both x86\_64 and ARM64; the execution environment is reproducible from source on any compatible host.

\subsection{Open Science}
\label{subsec:3.6.3}

The 50-day dataset (2 April -- 21 May 2026, Parquet format) is deposited on Zenodo~\cite{Zenodo} under CC~BY~4.0. Zenodo assigns a \ac{DOI} to the deposit, making the exact snapshot permanently citable. The pipeline source code --- collection, parsing, analysis scripts, and Dockerfile --- is on GitHub under the MIT licence, pinned to the version that produced the results in Chapter~\ref{ch:results}.
